\section{Preparación del entorno de trabajo}

En esta sección se describen los requisitos hardware y software necesarios para el desarrollo del proyecto, así como las herramientas utilizadas durante su implementación.

Se detallan las dependencias del entorno backend basado en Node.js y Express, la configuración de la base de datos PostgreSQL, el uso de contenedores Docker para la orquestación de servicios y la estrategia de desarrollo adoptada.

Asimismo, se documenta la configuración del sistema de control de versiones y el entorno de documentación empleado para la redacción de esta memoria.

\subsection{Requisitos hardware}

El desarrollo del proyecto se realizó en dos equipos con las siguientes características:

\begin{itemize}
    \item Sistema operativo: Windows 11 64 bits
    \item Procesador: Intel Core i5 8600K (Equipo de sobremesa) / Intel Core i5 11400H (Portátil)
    \item Memoria RAM: 16 GB
    \item Espacio en disco disponible: superior a 20 GB
\end{itemize}

No se requieren características hardware especiales más allá de las necesarias para ejecutar herramientas de desarrollo web modernas y contenedores Docker en caso de utilizar virtualización.

\subsection{Requisitos software}

Las herramientas utilizadas durante el desarrollo del proyecto han sido las siguientes:
\subsubsection*{Entorno de desarrollo}
\begin{itemize}
    \item Node.js v24.13.1
    \item npm 11.10.1 (incluido con Node.js)
    \item Next.js 16.1.6(framework de React para el desarrollo del frontend)
    \item React 19.x 
\end{itemize}
\subsubsection*{Herramientas adicionales}
\begin{itemize}
    \item Git 2.34.0.windows.1
    \item GitHub Desktop 3.5.4
    \item Visual Studio Code 1.109.5
    \item MiKTeX 26.1 (distribución \LaTeX)
    \item Strawberry Perl v5.42.0 (necesario para la ejecución de \texttt{latexmk})
\end{itemize}

% \subsubsection*{Entorno de contenedores}

% Para ejecutar el entorno en una configuración basada en contenedores se requiere:

% \begin{itemize}
%     \item Docker Desktop 4.62.0 (Windows 11, 64 bits)
%     \item Docker Compose (incluido en Docker Desktop)
% \end{itemize}

\subsection{Instalación del entorno local}
\subsubsection*{Node.js y npm}

Se instaló la versión LTS de Node.js desde su \href{https://nodejs.org}{sitio oficial}.
Tras la instalación, se verificó su correcta integración en el sistema mediante:

\begin{lstlisting}[language=bash]
node -v
npm -v
\end{lstlisting}

\subsubsection*{Creación del proyecto Next.js}

El proyecto se inicializó mediante el comando:

\begin{lstlisting}[language=bash]
npx create-next-app@latest nombre-proyecto
\end{lstlisting}

Este comando genera automáticamente la estructura del proyecto e instala las siguientes dependencias principales:

\begin{itemize}
    \item next
    \item react
    \item react-dom
\end{itemize}

Además, se configuran scripts de desarrollo y compilación en el archivo \texttt{package.json}.

Para iniciar el servidor de desarrollo, se ejecuta:
\begin{lstlisting}[language=bash]
npm run dev
\end{lstlisting}

El cual queda accesible en:

\begin{center}
\texttt{http://localhost:3000}
\end{center}


% \subsubsection*{PostgreSQL}

% PostgreSQL puede instalarse desde su sitio oficial:

% \begin{center}
%     \url{https://www.postgresql.org/download/windows/}
% \end{center}

% La instalación incluye el servidor de base de datos y la herramienta gráfica pgAdmin para la administración.

% Una vez instalado, puede verificarse su correcto funcionamiento mediante:

% \begin{lstlisting}[language=bash]
% psql --version
% \end{lstlisting}

% El servidor queda disponible por defecto en el puerto 5432.

\subsubsection*{Control de versiones}

El control de versiones se configuró mediante Git, inicializando el repositorio local con:

\begin{lstlisting}[language=bash]
git init
git add .
git commit -m "Initial commit"
\end{lstlisting}

Posteriormente, el repositorio fue vinculado a un repositorio remoto alojado en GitHub.

\subsubsection*{Entorno de documentación en \LaTeX}

La memoria se desarrolló en \LaTeX utilizando la distribución MiKTeX junto con Strawberry Perl, necesario para la ejecución de \texttt{latexmk}. La edición y compilación se realizaron en Visual Studio Code mediante la extensión \textit{LaTeX Workshop}.

La compilación del documento se realiza mediante:

\begin{lstlisting}[language=bash]
latexmk -pdf -shell-escape main.tex
\end{lstlisting}

También es posible compilarlo mediante la función integrada en el editor.
\paragraph{Configuración del PATH}

Durante la instalación fue necesario añadir manualmente al \texttt{PATH} del sistema las rutas correspondientes a MiKTeX y Strawberry Perl, ya que inicialmente el comando \texttt{latexmk} no era reconocido por la terminal. Tras actualizar la variable de entorno y reiniciar la sesión, el entorno quedó correctamente configurado.

% \subsection{Instalación del entorno dockerizado}

% El repositorio incluye un archivo \texttt{docker-compose.yml} en la raíz del proyecto que define los servicios necesarios para la ejecución del sistema en entorno de desarrollo.

% Desde la raíz del proyecto, la construcción y puesta en marcha de los contenedores se realiza mediante:

% \begin{lstlisting}[language=bash]
% docker compose up --build
% \end{lstlisting}

% Una vez iniciado, el sistema queda accesible en:

% \begin{itemize}
%     \item Backend (Express): \texttt{http://localhost:3000}
%     \item Base de datos PostgreSQL: puerto \texttt{5432}
% \end{itemize}

% Para detener los contenedores se utiliza:

% \begin{lstlisting}[language=bash]
% docker compose down
% \end{lstlisting}

% En caso de necesitar un reinicio completo del entorno (incluyendo la reinstalación de dependencias), se eliminan los volúmenes asociados mediante:

% \begin{lstlisting}[language=bash]
% docker compose down -v
% \end{lstlisting}

% Adicionalmente, el despliegue utiliza imágenes base oficiales de Node.js y PostgreSQL (disponibles en Docker Hub) para la construcción de los contenedores de backend y base de datos.

% \subsection{Estrategia de desarrollo adoptada}

% Durante el desarrollo del proyecto se adoptó una estrategia híbrida con el fin de optimizar los tiempos de compilación y recarga del frontend.

% Inicialmente, el sistema completo (frontend, backend y base de datos) se ejecutó mediante Docker Compose. Sin embargo, debido a la latencia introducida por el sistema de archivos compartido entre Windows y los contenedores Linux, se optó por ejecutar el frontend Angular directamente en entorno local mediante \texttt{ng serve}, manteniendo el backend y la base de datos en contenedores Docker.

% Esta configuración permite:
% \begin{itemize}
%     \item Recarga en caliente (Hot Reload) más rápida del frontend.
%     \item Aislamiento del backend y la base de datos.
%     \item Mayor agilidad en el desarrollo diario.
% \end{itemize}

% Actualmente, el servicio correspondiente al frontend se encuentra comentado en el archivo \texttt{docker-compose.yml}, manteniéndose activos únicamente los servicios de backend y base de datos.
